% -------------------------------------------------
% Updated Abstract and Keywords (replace existing sections)
% -------------------------------------------------
\begin{abstract}
Dense small-object detection is vital for public-health vector surveillance, agricultural phenotyping, and digital pathology, yet state-of-the-art detectors remain too heavyweight for edge deployment. This proposal introduces a systematic, multi-step compression framework that integrates architecture-aware structured pruning, knowledge distillation, and post-training quantisation across YOLO-family and RT-DETR detectors. The proposed methodology evaluates baseline detection performance across three specialized high-density datasets---\textit{Aedes aegypti} egg images, the CoNSeP histology set, and the Sorghum-Grain-Count benchmark. To systematically reduce computational complexity, the pipeline investigates channel-level structured pruning across individual and combined architectural modules (backbone, neck, and detection head), followed by knowledge distillation from uncompressed teacher models to recover degraded multi-scale representations, and uniform post-training quantisation ($\text{INT8}$) to minimize memory footprint. The study characterizes module vulnerability to parameter removal and evaluates the trade-offs between compression ratio, inference latency, and detection reliability (mean average precision and counting accuracy) to establish an empirical Pareto frontier. Ultimately, this research aims to deliver an architecture-aware compression pipeline and practical deployment guidelines for resource-constrained edge applications that require accurate dense counting.
\end{abstract}

\begin{flushleft}
\begin{tabular}{lp{4.25in}}
\textbf{Keywords:} & edge computing, small-object detection, model compression, knowledge distillation, quantisation \\
\end{tabular}
\end{flushleft}

